% ==================== 第2章 恒定电场 ====================
% 子文件：仅正文，由 1.电磁场笔记.tex \input 引入
% 教材：《工程电磁场导论》第2章 §2-1 ~ §2-5
% 公式统一用 formulaBox（本目录 3.矢量分析.tex 中定义）
% ==========================================================

\chapter{恒定电场}

\textbf{研究对象}：导电媒质中的恒定电场，即导电媒质内有恒定电流通过时，媒质内、外建立的电场。

\textbf{基本特征}：电荷分布不随时间变化，处处保持动态平衡；场量（$\vec{E}$、$\vec{J}$、$\rho$）均与时间无关，
\begin{formulaBox}
\[
\frac{\partial}{\partial t}(\bullet)=0
\]
\end{formulaBox}
这一条件把恒定电场与静电场、时变电磁场区分开：静电场由\emph{静止}电荷产生、场域内无电流；恒定电场由\emph{恒定运动}的电荷（恒定电流）维持，导电媒质内部有电流通过。因此恒定电场中\textbf{导体不再是等位体}，导体内部场强不为零（$E=J/\gamma \neq 0$），这是它与静电场最容易被混淆的区别。

本章按教材节序展开：导电媒质中的电流（\S2-1）$\to$ 电源电动势与局外场强（\S2-2）$\to$ 基本方程与分界面衔接条件（\S2-3）$\to$ 恒定电场与静电场的比拟（\S2-4）$\to$ 电导和部分电导（\S2-5）。

% ---------------------------------------------------------
\section{导电媒质中的电流}

\subsection{传导电流与运流电流}
按电流形成的物理机制分类：
\begin{itemize}
  \item \textbf{传导电流}：导电媒质（导体、电解液、大地等）中自由电荷的定向运动形成，服从欧姆定律；
  \item \textbf{运流电流}：自由空间（真空、气体）中带电体的机械运动形成，\textbf{不服从欧姆定律}。
\end{itemize}
本章只讨论传导电流。描述电流有两个层次的物理量：电流强度（整体量）与电流密度（分布量）。

\subsection{电流强度与电流密度矢量}
\textbf{电流强度}定义为单位时间内通过某一曲面的电量，
\begin{formulaBox}
\[
I=\frac{\mathrm{d}q}{\mathrm{d}t}\qquad(\mathrm{A})
\]
\end{formulaBox}
它是\textbf{通量}概念，只说明每秒通过某面积的总电量，无法反映电荷在截面各点的流动情况。

\textbf{电流密度矢量} $\vec{J}$ 描述电流的空间分布：方向取正电荷运动方向；大小为单位时间内穿过与该方向垂直的单位面积的净电量。对以速度 $\vec{v}$ 运动的体电荷分布，
\begin{formulaBox}
\[
\vec{J}=\rho\vec{v}\qquad(\mathrm{A/m^2})
\]
\end{formulaBox}
其中 $\rho$ 为电荷体密度。穿过任意曲面 $S$ 的电流是 $\vec{J}$ 的通量：
\begin{formulaBox}
\[
I=\iint_S \vec{J}\cdot\mathrm{d}\vec{S}
\]
\end{formulaBox}
按电荷分布形态，电流分为三类（$\vec{J}$、$\vec{K}$、$I$ 分别对应体、面、线分布）：
\begin{center}
\small
\begin{tabularx}{\linewidth}{@{}l l X@{}}
\toprule
\textbf{分布} & \textbf{密度} & \textbf{与电流强度的关系} \\
\midrule
体电流 & $\vec{J}=\rho\vec{v}$\ ($\mathrm{A/m^2}$) & $I=\displaystyle\iint_S \vec{J}\cdot\mathrm{d}\vec{S}$ \\
面电流 & $\vec{K}=\sigma\vec{v}$\ ($\mathrm{A/m}$) & $I=\displaystyle\int_l \vec{K}\cdot\vec{e}_n\,\mathrm{d}l$（$\vec{e}_n$ 为面法向） \\
线电流 & $I=\tau v$（$\tau$ 为单位长度电荷） & 沿导线流动，$I$ 即电流强度 \\
\bottomrule
\end{tabularx}
\end{center}

\corepoint{三层关系不要写反：$\vec{J}$ 是\emph{点}函数（$\mathrm{A/m^2}$），$\vec{K}$ 是\emph{面}上的密度（$\mathrm{A/m}$，沿面法向线元积分），$I$ 是\emph{积分}结果（$\mathrm{A}$）。凡是题目给"电流线密度 $K$"，就要用 $\int K\,\mathrm{d}l$ 而不是 $\iint J\,\mathrm{d}S$。}

\subsection{欧姆定律的微分形式}
在线性、各向同性导电媒质中，任一点传导电流密度与该点电场强度成正比：
\begin{formulaBox}
\[
\vec{J}=\gamma\vec{E}
\]
\end{formulaBox}
\begin{itemize}
  \item $\gamma$ 为媒质的\textbf{电导率}（$\mathrm{S/m}$），$\rho_r=1/\gamma$ 为电阻率（$\Omega\cdot\mathrm{m}$）；
  \item 适用条件：线性、各向同性导电媒质，且为\emph{点}的方程，与导体形状无关；
  \item 物理意义：导电媒质的物性方程，地位与电介质的 $\vec{D}=\varepsilon\vec{E}$、磁介质的 $\vec{B}=\mu\vec{H}$ 相当。
\end{itemize}
对一段均匀导体积分即得电路形式的欧姆定律 $U=IR$、$I=GU$，故 $U=IR$ 只是微分形式的宏观平均结果。

\subsection{焦耳定律的微分形式}
单位体积内的功率损耗（热功率体密度）为
\begin{formulaBox}
\[
p=\lim_{\Delta V\to 0}\frac{\Delta P}{\Delta V}=\vec{J}\cdot\vec{E}
\qquad(\mathrm{W/m^3})
\]
\end{formulaBox}
在线性、各向同性媒质中，代入 $\vec{J}=\gamma\vec{E}$ 得常用的三种等价写法：
\begin{formulaBox}
\[
p=\gamma E^2=\frac{J^2}{\gamma}
\]
\end{formulaBox}
对整个导体体积分得总功率损耗
\begin{formulaBox}
\[
P=\int_V p\,\mathrm{d}V=\int_V \vec{E}\cdot\vec{J}\,\mathrm{d}V=I^2R
\]
\end{formulaBox}
\textbf{物理意义}：恒定电场中电流的流动必然伴随能量损耗，要维持恒定电流必须持续由电源补充能量；$p$ 是\emph{点}的损耗，$P=I^2R$ 是整体结果，两者不能混用。

% ---------------------------------------------------------
\section{电源电动势与局外场强}

\subsection{局外力与局外场强}
要维持恒定电流，必须有非静电力把正电荷从低电位处"搬"回高电位处，这种力称为\textbf{局外力} $\vec{f}_e$（不是电磁力）。把作用于单位正电荷的局外力等效为一个场强，称为\textbf{局外场强}：
\begin{formulaBox}
\[
\vec{E}_e=\lim_{q\to 0}\frac{\vec{f}_e}{q}
\]
\end{formulaBox}
\begin{itemize}
  \item 方向：由电源负极指向正极（电源内部）；
  \item 它\textbf{不是}真正的电场强度，只在"引起电流"这一点上与电场强度等效；它不满足保守场条件；
  \item 电源外部不存在局外场强，只有库仑场强 $\vec{E}_c$。
\end{itemize}

\subsection{含源导电媒质中的电流与电动势}
电源内部总场强为 $\vec{E}_c+\vec{E}_e$，故通过含源导电媒质的电流密度为
\begin{formulaBox}
\[
\vec{J}=\gamma(\vec{E}_e+\vec{E}_c)
\]
\end{formulaBox}
\textbf{电动势}衡量电源做功能力，定义为单位正电荷从负极移到正极时局外力所做的功：
\begin{formulaBox}
\[
\varepsilon=\int_{B}^{A}\vec{E}_e\cdot\mathrm{d}\vec{l}
\]
\end{formulaBox}
（$B$ 为负极、$A$ 为正极，积分只在电源内部进行。）

\subsection{电源外恒定电场的守恒性}
对不经过电源的任意闭合路径 $L$，
\begin{formulaBox}
\[
\oint_L \vec{E}\cdot\mathrm{d}\vec{l}=\oint_L \vec{E}_c\cdot\mathrm{d}\vec{l}=0
\]
\end{formulaBox}
即\textbf{电源外恒定电场仍是无旋（保守）场}，可以引入标量电位 $\vec{E}=-\nabla\varphi$。

\corepoint{$\oint_L\vec{E}_e\cdot\mathrm{d}\vec{l}\neq 0$（局外场非保守），但 $\oint_L\vec{E}_c\cdot\mathrm{d}\vec{l}=0$。讨论"恒定电场是否无旋"时，指的是库仑场强 $\vec{E}_c$，且积分路径\emph{不穿过电源}；若路径穿过电源，总场强的环路积分等于电动势而不是零。}

% ---------------------------------------------------------
\section{恒定电场的基本方程与分界面上的衔接条件}

\subsection{电荷守恒定律与电流连续性方程}
单位时间流出闭合面 $S$ 的电荷等于体积 $V$ 内电量的减少率：
\begin{formulaBox}
\[
\oint_S \vec{J}\cdot\mathrm{d}\vec{S}=-\frac{\mathrm{d}q}{\mathrm{d}t}
\]
\end{formulaBox}
由高斯散度定理得\textbf{电流连续性方程}（微分形式）
\begin{formulaBox}
\[
\nabla\cdot\vec{J}+\frac{\partial\rho}{\partial t}=0
\]
\end{formulaBox}
恒定条件下 $\partial\rho/\partial t=0$，退化为
\begin{formulaBox}
\[
\nabla\cdot\vec{J}=0,\qquad \oint_S \vec{J}\cdot\mathrm{d}\vec{S}=0
\]
\end{formulaBox}
\textbf{物理意义}：恒定电流"无头无尾"，电流线是闭合曲线；任一闭合面净流出的传导电流为零（流入等于流出）。

\subsection{恒定电场的基本方程组}
电源外导电媒质中的恒定电场由下列方程完整描述（三种形式并列，缺一不可）：
\begin{formulaBox}
\[
\begin{cases}
\nabla\times\vec{E}=0 & \text{（无旋）}\\[2pt]
\nabla\cdot\vec{J}=0 & \text{（无源）}\\[2pt]
\vec{J}=\gamma\vec{E} & \text{（物性方程）}
\end{cases}
\qquad\Longleftrightarrow\qquad
\begin{cases}
\displaystyle\oint_L \vec{E}\cdot\mathrm{d}\vec{l}=0\\[6pt]
\displaystyle\oint_S \vec{J}\cdot\mathrm{d}\vec{S}=0
\end{cases}
\]
\end{formulaBox}
\begin{itemize}
  \item $\nabla\times\vec{E}=0$：恒定电场是无旋场，可引入电位 $\vec{E}=-\nabla\varphi$；
  \item $\nabla\cdot\vec{J}=0$：恒定电场是无源场\emph{（就电流线而言）}，但注意 $\nabla\cdot\vec{D}=\rho$ 仍可能成立（见下款）；
  \item 辅助方程（有电介质时）：$\nabla\cdot\vec{D}=\rho$，$\vec{D}=\varepsilon\vec{E}$，与静电场相同。
\end{itemize}

\subsection{分界面上的衔接条件}
设 $\vec{e}_n$ 为分界面上由媒质 1 指向媒质 2 的单位法向，由积分形式的基本方程可导出四个衔接条件（前两个是\emph{场量}形式，后两个是\emph{电位}形式）：
\begin{formulaBox}
\[
\begin{cases}
E_{1t}=E_{2t}\ \ \text{或}\ \ \vec{e}_n\times(\vec{E}_2-\vec{E}_1)=0 & \text{（切向：}E_t\text{连续）}\\[4pt]
J_{1n}=J_{2n}\ \ \text{或}\ \ \vec{e}_n\cdot(\vec{J}_2-\vec{J}_1)=0 & \text{（法向：}J_n\text{连续）}\\[4pt]
D_{2n}-D_{1n}=\sigma\ \ \text{或}\ \ \vec{e}_n\cdot(\vec{D}_2-\vec{D}_1)=\sigma & \text{（法向：}D_n\text{跃变）}\\[4pt]
\varphi_1=\varphi_2,\quad \gamma_1\dfrac{\partial\varphi_1}{\partial n}=\gamma_2\dfrac{\partial\varphi_2}{\partial n} & \text{（电位形式）}
\end{cases}
\]
\end{formulaBox}
\textbf{条件与含义}：$E_t$ 连续来自 $\oint\vec{E}\cdot\mathrm{d}\vec{l}=0$（沿跨界面小矩形回路）；$J_n$ 连续来自 $\oint\vec{J}\cdot\mathrm{d}\vec{S}=0$（跨界面小扁盒）；$D_n$ 的跃变量恰为分界面上的自由面电荷密度 $\sigma$。注意 $J_n$ 连续\textbf{不意味着} $D_n$ 连续——因为两侧的 $\varepsilon/\gamma$ 不同。

综合 $E_{1t}=E_{2t}$、$J_{1n}=J_{2n}$ 与 $\vec{J}=\gamma\vec{E}$，得\textbf{恒定电场的折射定律}：
\begin{formulaBox}
\[
\frac{\tan\theta_1}{\tan\theta_2}=\frac{\gamma_1}{\gamma_2}
\]
\end{formulaBox}
其中 $\theta_1,\theta_2$ 分别为两侧电流密度（或电场强度）与界面法线的夹角。

\subsection{两类特殊分界面的讨论}
\begin{enumerate}
  \item \textbf{良导体—不良导体}（$\gamma_2\gg\gamma_1$）：由折射定律 $\tan\theta_1=(\gamma_1/\gamma_2)\tan\theta_2\to 0$，故 $\theta_1\approx0$。结论：从不良导体一侧看，电流线（及 $\vec{E}$ 线）近似\emph{垂直于}分界面；良导体一侧的界面可近似当作\textbf{等位面}。
  \item \textbf{导体—理想介质}（如架空线与空气，$\gamma_2=0$）：由 $J_{2n}=J_{1n}$ 得 $J_{2n}=0$，即介质中无电流，导体侧电流只能沿表面流动（$J_1=J_{1t}$）；同时
  \[
  D_{2n}-D_{1n}=\sigma \Rightarrow D_{2n}=\sigma,\qquad E_{1t}=E_{2t},\ E_1=E_{1t}
  \]
  结论：分界面上\textbf{必有面电荷分布}；电场切向分量不为零，介质中紧贴导体表面的 $\vec{E}$ \emph{不再}垂直于导体表面。但通常 $E_{2n}\gg E_{2t}$，故工程上仍近似认为导体表面外的场垂直于表面。
\end{enumerate}

\corepoint{本节最易错的三处：(1) 把"$J_n$ 连续"写成"$D_n$ 连续"；(2) 把"$E_t$ 连续"与"$D_t$ 连续"混淆（$D_t$ 一般不连续）；(3) 直接套用静电场"导体表面场垂直于表面"的结论——在恒定电场中它只是近似。}

\subsection{导电媒质中恒定电荷的分布}
\begin{enumerate}
  \item \textbf{均匀媒质}（$\gamma$、$\varepsilon$ 均为常数）：由 $\nabla\cdot\vec{J}=0$ 得
  \begin{formulaBox}
  \[
  \nabla\cdot(\gamma\vec{E})=\nabla\cdot\Big(\frac{\gamma}{\varepsilon}\vec{D}\Big)=0
  \ \Longrightarrow\ \rho=0
  \]
  \end{formulaBox}
  即均匀导电媒质内部\textbf{无体电荷}，电荷只能分布在分界面上或电极表面。
  \item \textbf{不均匀媒质}（$\gamma$ 随位置变化）：
  \begin{formulaBox}
  \[
  \rho=\nabla\cdot\vec{D}=\nabla\cdot\Big(\frac{\varepsilon}{\gamma}\vec{J}\Big)
  =\vec{J}\cdot\nabla\Big(\frac{\varepsilon}{\gamma}\Big)
  \]
  \end{formulaBox}
  即媒质内部\textbf{可以存在体电荷}，其分布取决于 $\varepsilon/\gamma$ 的梯度方向与电流方向的关系。
  \item \textbf{分界面上的面电荷}：
  \begin{formulaBox}
  \[
  \sigma=D_{2n}-D_{1n}=\varepsilon_2E_{2n}-\varepsilon_1E_{1n}
  =\Big(\frac{\varepsilon_2}{\gamma_2}-\frac{\varepsilon_1}{\gamma_1}\Big)J_{1n}
  \]
  \end{formulaBox}
  对两种\emph{金属}导体的分界面，经典电子论给出各金属的 $\varepsilon$ 近似相同，故 $\sigma\approx0$；而对金属—介质、不同电解液的分界面一般 $\sigma\neq0$。
\end{enumerate}

\subsection{恒定电场的边值问题}
由 $\nabla\times\vec{E}=0$ 引入 $\vec{E}=-\nabla\varphi$，代入 $\nabla\cdot(\gamma\vec{E})=0$。均匀媒质中 $\gamma$ 为常数，得
\begin{formulaBox}
\[
\nabla^2\varphi=0
\]
\end{formulaBox}
即恒定电场的电位满足\textbf{拉普拉斯方程}，配合衔接条件 $\varphi_1=\varphi_2$、$\gamma_1\partial\varphi_1/\partial n=\gamma_2\partial\varphi_2/\partial n$ 与场域边界条件（给定电位的第一类边界、给定法向电流密度的第二类边界），构成恒定电场的\textbf{边值问题}。它与静电场边值问题同型，故解法（分离变量、镜像、有限差分）完全可借用——这正是下一节"比拟"的依据。

% ---------------------------------------------------------
\section{导电媒质中的恒定电场与静电场的比拟}

\subsection{比拟的依据与条件}
比较\emph{电源外}导电媒质中的恒定电场与\emph{无电荷分布区}的静电场：
\[
\text{静电场：}\ \nabla\times\vec{E}=0,\ \nabla\cdot\vec{D}=0
\qquad
\text{恒定电场：}\ \nabla\times\vec{E}=0,\ \nabla\cdot\vec{J}=0
\]
两者的\emph{无旋}方程完全相同，\emph{无源}方程形式相同，本构方程 $\vec{D}=\varepsilon\vec{E}$ 与 $\vec{J}=\gamma\vec{E}$ 一一对应，衔接条件也一一对应。因此，只要两个场的\textbf{边界条件相同}，其电位函数 $\varphi$ 的解答就完全相同，场图（等位面、场线）一致。

\textbf{适用条件}（缺一不可）：
\begin{enumerate}
  \item 恒定电场须取\emph{电源外}区域（否则 $\oint\vec{E}\cdot\mathrm{d}\vec{l}\neq0$，无旋性不成立）；
  \item 静电场须取\emph{无自由电荷}区域（$\rho=0$），否则 $\nabla\cdot\vec{D}\neq0$；
  \item 若媒质分区不均匀，还须满足 $\varepsilon_1/\varepsilon_2=\gamma_1/\gamma_2$（即各分区内 $\varepsilon/\gamma$ 为同一常数），否则分界面折射情况不同，比拟不成立。
\end{enumerate}

\subsection{对应量对照表}
\begin{center}
\small
\begin{tabularx}{\linewidth}{@{}l X X@{}}
\toprule
\textbf{物理量/关系} & \textbf{静电场（$\rho=0$）} & \textbf{恒定电场（电源外）} \\
\midrule
基本方程（无旋） & $\nabla\times\vec{E}=0$，$\oint\vec{E}\cdot\mathrm{d}\vec{l}=0$ & 同左 \\
基本方程（无源） & $\nabla\cdot\vec{D}=0$，$\oint\vec{D}\cdot\mathrm{d}\vec{S}=0$ & $\nabla\cdot\vec{J}=0$，$\oint\vec{J}\cdot\mathrm{d}\vec{S}=0$ \\
本构方程 & $\vec{D}=\varepsilon\vec{E}$ & $\vec{J}=\gamma\vec{E}$ \\
场量 & $\vec{E}$，$\vec{D}$ & $\vec{E}$，$\vec{J}$ \\
位函数 & 电位 $\varphi$，$\vec{E}=-\nabla\varphi$ & 电位 $\varphi$，$\vec{E}=-\nabla\varphi$ \\
控制方程 & $\nabla^2\varphi=0$ & $\nabla^2\varphi=0$ \\
衔接条件 & $\varphi_1=\varphi_2$，$\varepsilon_1\partial\varphi_1/\partial n=\varepsilon_2\partial\varphi_2/\partial n$ & $\varphi_1=\varphi_2$，$\gamma_1\partial\varphi_1/\partial n=\gamma_2\partial\varphi_2/\partial n$ \\
源量 & 电荷 $q$，面电荷 $\sigma$ & 电流 $I$，面电流 $K$ \\
集总参数 & 电容 $C=q/U$ & 电导 $G=I/U$ \\
参数关系 & $C=\dfrac{\varepsilon\int_S\vec{E}\cdot\mathrm{d}\vec{S}}{\int_A^B\vec{E}\cdot\mathrm{d}\vec{l}}$ & $G=\dfrac{\gamma\int_S\vec{E}\cdot\mathrm{d}\vec{S}}{\int_A^B\vec{E}\cdot\mathrm{d}\vec{l}}$ \\
能量 & 静电能量 $W=\frac12\int\vec{D}\cdot\vec{E}\,\mathrm{d}V$ & 功率损耗 $P=\int\vec{J}\cdot\vec{E}\,\mathrm{d}V$ \\
\bottomrule
\end{tabularx}
\end{center}
由表中两式相比立得核心关系式
\begin{formulaBox}
\[
\frac{C}{G}=\frac{\varepsilon}{\gamma}
\qquad\Longleftrightarrow\qquad
G=\frac{\gamma}{\varepsilon}C
\]
\end{formulaBox}
\textbf{用法}：已知（或易求）电容 $C$，把 $\varepsilon\to\gamma$ 即得电导 $G$；反之亦然。这是计算电导最省力的途径。

\subsection{比拟法的三种用法}
\begin{enumerate}
  \item \textbf{参数置换}：由已知的电容公式直接写出电导公式（$C\to G$，$\varepsilon\to\gamma$）；
  \item \textbf{镜像法}：导电媒质分界面（如土壤—空气）中的线电流问题，可用静电场中点电荷对介质分界面的镜像结果换算（$\varepsilon\to\gamma$，$q\to I$）；
  \item \textbf{模拟实验}：用导电溶液（电解槽）中的恒定电流场模拟难以直接测量的静电场（如高压套管、电缆头、绝缘子电场），测电位分布即得原静电场分布，称为\textbf{电解槽模拟}。
\end{enumerate}

\corepoint{比拟是"同型方程 + 同边界条件"的数学对应，不是物理等价。使用时必须声明三条适用条件，尤其是"电源外"和"$\varepsilon/\gamma$ 为常数"。若媒质不均匀且不满足 $\varepsilon_1/\varepsilon_2=\gamma_1/\gamma_2$，只能逐区求解，不能整体置换。}

% ---------------------------------------------------------
\section{电导和部分电导}

\subsection{电导与电阻的定义}
\textbf{电导} $G$ 定义为流经导电媒质的电流与两端电压之比（电阻 $R=1/G$）：
\begin{formulaBox}
\[
G=\frac{I}{U}
=\frac{\displaystyle\iint_S \vec{J}\cdot\mathrm{d}\vec{S}}
     {\displaystyle\int_A^B \vec{E}\cdot\mathrm{d}\vec{l}}
=\frac{\gamma\displaystyle\iint_S \vec{E}\cdot\mathrm{d}\vec{S}}
     {\displaystyle\int_A^B \vec{E}\cdot\mathrm{d}\vec{l}}
\qquad(\mathrm{S})
\]
\end{formulaBox}
可见 $G$ 只与两电极的形状、尺寸、相互位置及媒质电导率有关，与所加电压、电流大小无关；\emph{电阻（电导）的计算实质上是电场的计算}。

\textbf{三条计算路线}：
\begin{enumerate}
  \item 假设电流 $I$：$I\to\vec{J}\to\vec{E}=\vec{J}/\gamma\to U=\int\vec{E}\cdot\mathrm{d}\vec{l}\to G=I/U$（有对称性时最常用）；
  \item 假设电压 $U$：$U\to\vec{E}\to\vec{J}=\gamma\vec{E}\to I=\iint\vec{J}\cdot\mathrm{d}\vec{S}\to G$（不常用）；
  \item 解拉普拉斯方程求 $\varphi$，再积分求 $I$、$U$（一般情形）；
  \item 静电比拟：先求同形电极系统的 $C$，再由 $G=(\gamma/\varepsilon)C$ 换算。
\end{enumerate}

\textbf{三个必须记住的结果}（均可由比拟法立即得到）：
\begin{formulaBox}
\[
G_{\text{平板}}=\frac{\gamma S}{d},\qquad
G_{\text{同轴}}=\frac{2\pi\gamma l}{\ln(R_2/R_1)},\qquad
G_{\text{同心球}}=\frac{4\pi\gamma R_1R_2}{R_2-R_1}
\]
\end{formulaBox}
同轴电缆的\textbf{绝缘电阻}即 $R=1/G=\dfrac{\ln(R_2/R_1)}{2\pi\gamma l}$，它是用绝缘材料隔开的两个导体之间的直流电阻，由绝缘缺陷（气隙、开裂、受潮、老化）引起的微弱泄漏电流决定，可用于判断绝缘状态。

\subsection{部分电导}
当导电媒质中有三个及以上良导体电极（可视为等位体）时，任意两电极间的电流不仅取决于它们自身的电压，还受其它电极电压影响，单一电导不足以描述，需引入\textbf{部分电导}。

设 $(n+1)$ 个电极（0 号取作参考电极），线性各向同性媒质中由叠加原理可写出各电极电流与电压的关系：
\begin{formulaBox}
\[
\begin{cases}
I_1=G_{10}U_{10}+G_{12}U_{12}+\cdots+G_{1n}U_{1n}\\
I_2=G_{21}U_{21}+G_{20}U_{20}+\cdots+G_{2n}U_{2n}\\
\qquad\vdots
\end{cases}
\qquad U_{kj}=\varphi_k-\varphi_j
\]
\end{formulaBox}
\begin{itemize}
  \item $G_{k0}$ 称为\textbf{自有部分电导}（电极 $k$ 与参考电极之间）；
  \item $G_{kj}$（$k\neq j$）称为\textbf{互有部分电导}（电极 $k$、$j$ 之间）；
  \item 全部部分电导均为正值（互有部分电导描述两电极间的直接漏电通道，故与"互有电导"式的负系数不同）；
  \item $(n+1)$ 个电极的系统共有 $\dfrac{n(n+1)}{2}$ 个部分电导，与部分电容一一对应（$C_{kj}\leftrightarrow G_{kj}$，$\varepsilon\leftrightarrow\gamma$）。
\end{itemize}
工程上"三极法"测接地电阻、多电极系统的漏电分析都用部分电导描述。

\subsection{接地电阻}
把电气设备的一部分与大地连接称为\textbf{接地}；埋入地中并与设备相连的导体（系统）称为\textbf{接地体}，接地体与接地线合称接地装置。\textbf{接地电阻}是电流由接地装置流入大地、再经大地向远处扩散所遇到的电阻，包括接地线与接地体本身电阻、接地体与大地间的接触电阻、大地电阻，其中\textbf{以前者为主的是大地电阻}（前三项远小于最后一项）。

计算时把接地体看作电极、取无穷远为零电位，且因电流在电极附近密度最大，接地电阻主要集中于电极附近，故远端电流如何汇聚影响很小。

\textbf{深埋导体球}（半径 $a$，可不计地面影响）：$J=I/(4\pi r^2)$，$E=I/(4\pi\gamma r^2)$，
\begin{formulaBox}
\[
R=\int_a^\infty \frac{I}{4\pi\gamma r^2}\,\mathrm{d}r=\frac{1}{4\pi\gamma a}
\]
\end{formulaBox}
\textbf{紧靠地面的半球形接地体}：用镜像法，实际电极与其镜像构成完整球，系统中流出的电流为实际的两倍，故实际接地电阻为完整球的两倍：
\begin{formulaBox}
\[
R=2\cdot\frac{1}{4\pi\gamma a}=\frac{1}{2\pi\gamma a}
\]
\end{formulaBox}

\subsection{跨步电压与危险区}
接地体附近有大的入地电流时，地面上相距一个步距的两点间存在电位差，称为\textbf{跨步电压}。对半径为 $a$ 、入地电流为 $I$ 的半球形接地体，距球心 $r$ 处
\begin{formulaBox}
\[
J=\frac{I}{2\pi r^2},\qquad E=\frac{I}{2\pi\gamma r^2},\qquad
\varphi(r)=\frac{I}{2\pi\gamma r}
\]
\end{formulaBox}
设 $A$、$B$ 两点与球心距离分别为 $l$、$l+b$（$b$ 为步距），则跨步电压
\begin{formulaBox}
\[
U_{BA}=\varphi_A-\varphi_B=\frac{I}{2\pi\gamma}\Big(\frac{1}{l}-\frac{1}{l+b}\Big)
=\frac{I\,b}{2\pi\gamma\,l(l+b)}
\]
\end{formulaBox}
若对人体有危险的临界电压为 $U_0$，令 $U_{BA}=U_0$ 解出的 $l$ 即\textbf{危险区}边界，危险区是以接地体中心为圆心、半径 $l$ 的圆面积。当 $l\gg b$ 时可近似
\begin{formulaBox}
\[
U_{BA}\approx\frac{I\,b}{2\pi\gamma l^2}
\ \Longrightarrow\
l\approx\sqrt{\frac{I\,b}{2\pi\gamma U_0}}
\]
\end{formulaBox}
真正危及生命的不是电压而是通过人体的电流（工频电流超过约 30\,mA 即有生命危险），故 $U_0$ 应按人体允许电流与人体电阻折算。

\corepoint{接地电阻常见三个错误：把"接地体本身的电阻"当接地电阻（它只占极小部分）；忘记半球形接地体要\emph{乘 2}（镜像法带来的倍数）；把跨步电压写成 $\varphi_A+\varphi_B$ 或漏掉步距 $b$ 的影响。}
